\documentclass{article}
\usepackage{amsmath, amssymb, amsthm}
\usepackage{hyperref}
\usepackage{geometry}
\geometry{margin=1in}

\title{Erd\H{o}s Problem \#66}
\date{Last updated: 2025-08-31}
\author{Source: erdosproblems.com}

\begin{document}
\maketitle

\noindent\textbf{Status:} OPEN \quad \textbf{Prize: \$500}

\noindent\textbf{Tags:} number theory, additive basis

\medskip
\noindent\textbf{Problem Statement:}

Is there $A\subseteq \mathbb{N}$ such that\[\lim_{n\to \infty}\frac{1_A\ast 1_A(n)}{\log n}\]exists and is $\neq 0$?

\bigskip
\noindent\textbf{Remarks:}

A suitably constructed random set has this property if we are allowed to ignore an exceptional set of density zero. The challenge is obtaining this with no exceptional set. Erd\H{o}s believed the answer should be no. In \cite{Er80} he explicitly asked whether there exists such a set where the limit is $1$ (and did not believe there existed such a sequence).

Erd\H{o}s and S\'{a}rk\"{o}zy proved that\[\frac{\lvert 1_A\ast 1_A(n)-\log n\rvert}{\sqrt{\log n}}\to 0\]is impossible. Erd\H{o}s suggests it may even be true that the $\liminf$ and $\limsup$ of $1_A\ast 1_A(n)/\log n$ are always separated by some absolute constant. 

Horv\'{a}th \cite{Ho07} proved that\[\lvert 1_A\ast 1_A(n)-\log n\rvert \leq (1-\epsilon)\sqrt{\log n}\]cannot hold for all large $n$.

\bigskip
\noindent\textbf{References:}

[Er80] Erd\H{o}s, Paul, A survey of problems in combinatorial number theory. Ann. Discrete Math. (1980), 89-115.
 
 [Ho07] G. Horv\'{a}th, An improvement of a theorem of Erd\H{o}s and S\'{a}rk\"{o}zy. Pollack Periodica (2007), 155-161.

\bigskip
\noindent\small{Source: \url{https://www.erdosproblems.com/66}}
\end{document}
